\chapter{DNA-aware Transformer mechanisms}
\label{evod24:chapter}
Chapter 23 made Hermon's Transformer algebra explicit. That does not yet make an
architecture DNA-aware. DNA-specific information has to enter through mechanisms
whose effect can be stated, implemented and tested: base spans, strand actions,
physical coordinate distances and legal long-context access.

Consider a variant changing one base in a six-base record. Which tokens change?
Does a position mean a token index or a base coordinate? If the record is reverse
complemented, which features should change sign, which positions should reverse,
and which label should remain the same? Finally, can distant positions communicate
without allocating a dense attention matrix? This chapter answers those questions
with small differentiable references, not a new untested architectural name.

\section{A token must carry its genomic span}
Let a canonical sequence be $x=(x_0,\ldots,x_{L-1})\in\{A,C,G,T\}^L$, $L>0$.
A token record is
\[
t_i=(u_i,a_i,b_i),\qquad 0\le a_i<b_i\le L,
\]
where $u_i=x[a_i:b_i]$ and $[a_i,b_i)$ is a zero-based half-open base span.
This local window convention is not automatically the coordinate system of a
reference assembly; that mapping must also name the assembly, chromosome and
orientation. Chapter 25 consolidates those dataset-level contracts.

For width $k$ and stride $s$, the executable reference emits spans
$[js,js+k)$ while the end stays inside the record. It rejects unsupported symbols,
empty records and invalid widths/strides rather than silently converting ambiguity
to a canonical base. Ambiguity handling is a distinct tokenizer policy.
\Listing{tokenize}
\begin{center}\small\input{\Generated/spans.tex}\end{center}
The generated table uses ACGTAC, $k=3$, $s=1$. A substitution at base coordinate
two changes tokens zero, one and two. Token three does not include that base.
The exact mutation footprint is
\[
\mathcal F(p)=\{i:a_i\le p<b_i\}.
\]
This is a coordinate statement, not a claim that all affected model activations
remain local after attention.

\Plate{01-spans}{One substitution changes every overlapping token whose span
contains the altered base. The base axis and token spans remain visible, so the
number of changed tokens is not confused with the number of changed nucleotides.}
{evod24:spans}

Overlapping $k$-mers do not necessarily shorten token count: at stride one,
$T=L-k+1$. Larger stride can shorten the sequence but leave terminal bases
unrepresented. A learned variable-span tokenizer must expose an analogous mapping
if variant attribution and physical distances are to remain meaningful.
DNABERT and DNABERT-2 supply established tokenizer design points, not evidence
that Hermon's chosen tokenizer is optimal \cite{evod24:dnabert,evod24:dnabert2}.

\subsection{Reverse complement acts on spans, not only strings}
A local span transforms under reverse complement as
\begin{equation}
(u,a,b)\longmapsto
\bigl(R(u),L-b,L-a\bigr).
\label{evod24:spanrc}
\end{equation}
Applying this twice returns the original span. Reversing a list of all stride-one
$k$-mers after this transformation equals tokenizing $R(x)$ with the same policy.

That equality can fail for strided tokenization. For ACGTA, $k=2$, stride two,
forward spans are $[0,2)$ and $[2,4)$, leaving base four uncovered. Their reversed
spans are $[1,3)$ and $[3,5)$, whereas independently tokenizing the RC record
starts again at zero. The tokenizer-grid origin changed.
RC consistency requires a boundary/grid policy; complementing token IDs alone
does not repair the missed-base geometry.

\section{Token causality is weaker than base causality}
Suppose the prediction target is base $x_p$ and legal context is $x[0:p]$.
An input token containing any base at coordinate $p$ or later is illegal,
even if its token index lies below the diagonal in a token-attention matrix.

For example, token ACG spans $[0,3)$ and already contains target G at base two.
Using that token to predict the third base is leakage. A triangular attention
mask cannot remove information encoded inside the input token itself.
The necessary span-availability condition is
\begin{equation}
b_i\le p
\quad\text{for every feature used to predict }x_p.
\label{evod24:availability}
\end{equation}
\Listing{causal_span_leaks}

\Plate{02-access}{The information boundary is a base coordinate. A token ending
after that boundary is illegal even when the token-level attention edge is causal.
No full-record RC, corruption cache or pooled suffix feature may bypass this gate.}
{evod24:access}

For causal token prediction, the prediction target can instead be the next token
identity. Then overlapping adjacent tokens may reveal most of that next identity;
the conditional entropy and interpretation differ from next-base prediction.
Declare which target is scored. A low overlapping-token loss is not automatically
a comparable nucleotide likelihood.

Masked-token objectives need a related check. Masking one token can leave its
target bases visible inside neighboring overlapping tokens. To obscure base $p$,
hide or consistently corrupt all spans in $\mathcal F(p)$ and any derived cached
features carrying the same information. This chapter computes the footprint; it
does not implement a masked pretraining pipeline. Causal, masked and whole-record
objectives remain distinct, and Chapter 21's offline dual-state representation
must not silently feed the causal case.

\section{Build a strand action into a feature space}
An unrestricted embedding lookup gives each base an unrelated learned vector.
It does not impose complement symmetry. A small alternative makes the symmetry
inspectable by separating even and odd feature channels.

Define pair class $g(A)=g(T)=0$, $g(C)=g(G)=1$ and signs
$\sigma(A)=\sigma(C)=+1$, $\sigma(T)=\sigma(G)=-1$.
Let $E_g,O_g,s\in\mathbb R^D$ be learned vectors, and let record orientation
$o\in\{-1,+1\}$ be declared metadata, not a target-derived label.
At base $i$, define
\begin{equation}
F_i(x,o)=
\begin{bmatrix}
E_{g(x_i)}\\
\sigma(x_i)O_{g(x_i)}+os
\end{bmatrix}
\in\mathbb R^{2\times D}.
\label{evod24:strand-definition}
\end{equation}
Complement preserves pair class and reverses sign. Define the joint record action
$(x,o)\mapsto(R(x),-o)$ and feature action
$(JF)_i=(F_{L-1-i}^{\mathrm{even}},-F_{L-1-i}^{\mathrm{odd}})$.
Then the construction satisfies
\[
F(R(x),-o)=JF(x,o).
\]
This is an exact algebraic equivariance contract. It is not a representation of a
physical DNA strand's electronic structure or a newly discovered biological law.
\Listing{StrandEmbedding.forward}

\Plate{03-strand}{Complement-paired bases share even features and reverse odd
features. Record reversal permutes positions; orientation metadata changes sign.
The symmetry is imposed by parameter tying, not inferred from the colors.}
{evod24:strand}

\begin{center}\small\input{\Generated/strand_features.tex}\end{center}
The generated trace assigns $E_0=(1,0)$, $E_1=(0,1)$,
$O_0=(.5,.25)$, $O_1=(-.25,.5)$ and $s=(.1,.2)$ with $o=+1$.
For A the odd channel is $(.6,.45)$; for its complemented T in orientation $-1$
it is $(-.6,-.45)$. Every value is generated from the same reference that tests
the joint action.

\subsection{Equivariant features do not make the whole model equivariant}
Flattening the two channels and applying an arbitrary linear map can mix even
and odd components in a way that does not commute with $J$.
An unrestricted Transformer can therefore break the property despite a correct
embedding. A pointwise linear map that commutes with the sign action must preserve
the parity subspaces; nonlinearities and attention require their own action
contracts. Ordinary additive positional embeddings can also break reversal
symmetry.

Chapter 23's shared two-pass offline readout provides a simple invariant baseline:
average the same function on a record and its joint RC transform.
By contrast, using different learned strand embeddings in two branches and merely
averaging does not guarantee invariance. Swapping the records does not swap
identical functions unless the parameter/metadata action has been specified.

Caduceus establishes RC-equivariant, bidirectional sequence modeling as prior art
\cite{evod24:caduceus}. Our tied feature reference is a small teaching construction,
not a reproduction of its architecture or evidence of improved genomic accuracy.
For orientation-sensitive labels, invariant pooling can discard the signal the
task requires. Use equivariant outputs or an explicit label action instead.

\section{Position should measure the distance the claim uses}
Chapter 23 used ordinal positions of valid bases to make inserted padding inert.
For variable-span tokens, token ordinal is not base distance. An optional
representative coordinate is a span's base-center midpoint,
\[
p_i=(a_i+b_i-1)/2.
\]
Even-width spans can have half-integer centers. A distance-sensitive feature should
state whether its positions are base coordinates, normalized window coordinates,
token ordinals or a multi-scale combination.

For example, tokens spanning $[0,2)$ and $[2,10)$ have centers $.5$ and $5.5$:
their base distance is five, not the ordinal distance one. Compressing unknown
bases or omitted regions changes those distances unless the original coordinate
mapping is retained. Padding is not a biological deletion.

\section{Derive rotary positions as pair rotations}
Rotary position embedding is an established Transformer mechanism
\cite{evod24:rope}. We derive its small reference to make the unit contract clear.
For a pair of coordinates and angular frequency $\omega$ in radians per base,
let
\[
\mathsf R(\theta)=
\begin{pmatrix}\cos\theta&-\sin\theta\\\sin\theta&\cos\theta\end{pmatrix},
\quad
\tilde q_i=\mathsf R(\omega p_i)q_i,\quad
\tilde k_j=\mathsf R(\omega p_j)k_j.
\]
Orthogonality and the angle-addition law give
\begin{equation}
\tilde q_i^\top\tilde k_j
=q_i^\top\mathsf R\!\left(\omega(p_j-p_i)\right)k_j.
\label{evod24:rope}
\end{equation}
Thus simultaneous translation $p_i\mapsto p_i+a$ cancels in pair scores.
The norm of each vector is unchanged. These claims hold pair by pair in an even
head dimension $D_h$, using $D_h/2$ frequencies.

\Plate{04-rotary}{Assigned two-dimensional query/key vectors are rotated by their
declared base positions. The dot product depends on relative angle; a common
coordinate translation cancels. This geometric identity is not a guarantee of
long-context generalization.}{evod24:rotary}
\Listing{rotary}

For $q=k=(1,0)$, positions zero and two, and $\omega=\pi/4$ rad/base, the rotated
vectors are $(1,0)$ and $(0,1)$ up to numerical error. Their score is zero.
Translating both positions by ten changes the individual directions but not their
dot product. The tests check this assigned case, norm preservation, explicit
block-diagonal matrix parity and input gradients in float64.

Relative-position algebra does not establish RC equivariance. Reversal sends
$p_j-p_i$ to its negative; content features also transform. Without a compatible
feature/score action, Equation~\ref{evod24:rope} can change. Similarly, a rotation
formula accepting large coordinates does not prove that a trained model
extrapolates to those coordinates. High-frequency phase aliasing and unseen
distance distributions require evaluation.

A multi-scale variant could place one rotary block on base distance and another
on larger genomic bins. Bin size, origin and boundary behavior are then model
assumptions. Such a variant is a research extension here, not an implemented
accuracy gain. An ablation must compare it against plain base positions under
matched records, objective and tuning budgets.

\section{Long context is an attention graph before it is a speed claim}
Let $\mathcal N(i)$ be the legal key set for query $i$. For local radius $w$,
a declared static global-position set $G$ and valid-position indicators $v_i$,
an offline edge exists when
\[
v_iv_j=1
\quad\text{and}\quad
\bigl(|i-j|\le w\ \lor\ i\in G\ \lor\ j\in G\bigr).
\]
The causal case additionally requires $j\le i$.
Global edges do not override causality. Padded query rows and padded keys remain
blocked as in Chapter 23.
\Listing{edge_lists}

Local-plus-global attention is established prior art, including Longformer
\cite{evod24:longformer}. Our graph is a transparent reference, not a novel
mechanism merely because its input is DNA.
At fixed $w$ and $g=|G|$, an edge-count upper bound is
\[
|\mathcal E|\le T(2w+1+2g).
\]
One $Tg$ term accounts for global keys available to local queries; the other
accounts for global queries attending broadly. Intersections and boundaries
reduce the actual count.
\begin{center}\small\input{\Generated/cost.tex}\end{center}
The table counts actual edges at radius four with the first and last positions
global. It does not report measured latency, FLOPs or peak memory.

\Plate{05-relay}{A distant source reaches a local query through a global relay in
two layers, not one. At layer one the global query gathers; at layer two another
query reads the updated global representation. Causal graphs cannot relay
future information backward.}{evod24:relay}

\subsection{Dependency paths have a layer count}
Let $\mathcal D_0(i)=\{i\}$ because a residual path preserves self-information.
With the same graph in each layer,
\[
\mathcal D_{\ell+1}(i)
=\mathcal D_\ell(i)\cup
 \bigcup_{j\in\mathcal D_\ell(i)}\mathcal N(j).
\]
In a length-twelve offline record with radius one, query one reaches positions
zero through three after two local layers. With position zero global, query one
still cannot access source eleven through that global node in the first layer:
the global node's update is being computed concurrently. It can access source
eleven after the second layer. The reference expands the dependency sets exactly.
\Listing{reachable_sources}

At causal query one, the same global policy cannot expose any source above one,
regardless of layer count. This is a pathwise proof that all edges respecting
causality preserve causality; the finite suffix-intervention test checks the
actual numerical implementation as well.

Global positions selected by inspecting the full sequence can create a different
leak even if edges look triangular. A future motif could determine which prefix
features are designated global. The causal policy must depend only on legal
metadata or available prefix information. Known full-record length and a streaming
endpoint are also different interface assumptions, addressed in Chapter 26.

\section{Implement sparse semantics without allocating dense scores}
A dense $T\times T$ score matrix filled with masked values is semantically sparse
but still consumes dense score storage. The reference instead stores ragged legal
key lists. For each query, it gathers only allowed keys and values, computes the
row scores, normalizes over those keys and reduces their values:
\[
a_{ij}=\frac{\exp(q_i^\top k_j/\sqrt{D_h})}
{\sum_{j'\in\mathcal N(i)}\exp(q_i^\top k_{j'}/\sqrt{D_h})},
\qquad
o_i=\sum_{j\in\mathcal N(i)}a_{ij}v_j.
\]
An empty row returns zero under the declared padding convention.
\Listing{sparse_attention}

\Plate{06-kernel}{Ragged gather processes only declared edges; the independent
dense oracle materializes a small full score matrix for value and gradient checks.
Semantic parity does not imply that a Python loop is a fast production kernel.}
{evod24:kernel}

For fixed $D_h$, score and weighted-value work scales with the number of edges.
Edge construction, sorting, projection work, autograd buffers and output storage
remain additional costs. The Python implementation is a pedagogical single-head
reference: launches and irregular gathers can make it slower than an optimized
dense kernel at modest lengths. We report no speedup.

A dense vectorized oracle independently reconstructs legal rows on small tensors.
Tests compare output values and all Q/K/V gradients in float64, including a
padding row. Finite differences check sparse derivatives. A future suffix
intervention changes offline prefix outputs but leaves causal prefix outputs
exactly unchanged under the tested graph. These are software contracts, not
evidence that the model learned regulatory grammar.

\section{Retrieval adds information, not free context}
A retrieval mechanism can supplement local context with selected external
sequences or annotations. It must declare the corpus version, coordinate mapping,
selection rule and permitted information. Retrieving a test label, a homologous
training duplicate or a suffix-derived feature changes the evaluation contract.
Static external knowledge and dynamic within-record access are not the same.

For causal prediction, retrieval queries and retrieved content must respect the
available-prefix and knowledge-cutoff assumptions. For offline variant scoring,
paired reference/alternate retrieval must not receive different target-derived
evidence. Retrieval is not implemented by the local-global graph above: a global
token summarizes the current record, not an external database. We keep retrieval
as a clearly labeled extension rather than calling ordinary attention external
memory.

\section{Design the experiment that could reject a mechanism}
The executable chapter establishes span, strand, rotary and sparse-graph
correctness. It does not train a new full genomic model.
A serious evaluation now asks separate questions:
\begin{enumerate}
\item Does a span-aware tokenizer improve variant-local attribution relative to
the same model with ordinal-only positions?
\item Does the declared RC action help invariant labels without damaging
orientation-sensitive tasks?
\item Does base-coordinate rotary position encoding improve the chosen context
distribution relative to Chapter 23's position baseline?
\item Does sparse access preserve task accuracy while meeting a declared
memory/latency budget?
\end{enumerate}
Each comparison freezes dataset, split groups, objective, width/depth budget,
training tokens, tuning effort and evaluation examples. Compare conventional
Transformer, local CNN and an appropriate non-attention long-context baseline,
not only DOGMA versus Hermon.

A broken mutation footprint falsifies coordinate correctness. A nonzero joint-RC
equivariance error falsifies the tied-feature construction. Gradient disagreement
falsifies implementation parity. Matched-budget failure to improve the task
weakens a usefulness claim even when every algebraic test passes.
Chapter 23's locally rerun three-seed pilot remains in its own evidence record;
the incoming package's old single-run percentage is provenance, not a current
summary silently carried into this chapter.

\section{Exercises and worked solutions}
\begin{enumerate}
\item List width-three spans of ACGTAC.
\textbf{Answer:} ACG $[0,3)$, CGT $[1,4)$, GTA $[2,5)$ and TAC $[3,6)$.
\item Which tokens change if base two changes?
\textbf{Answer:} zero, one and two, because exactly those spans contain two.
\item Derive Equation~\ref{evod24:spanrc}.
\textbf{Answer:} reversal sends base $p$ to $L-1-p$, mapping
$a,\ldots,b-1$ to $L-b,\ldots,L-a-1$; complement transforms the substring.
\item Why can a strided grid break RC commutation?
\textbf{Answer:} uncovered terminal bases move to the other end, changing the
grid origin; the ACGTA width-two stride-two example gives different spans.
\item Why can causal attention on ACG leak target G?
\textbf{Answer:} target base two is already inside span $[0,3)$.
\item To obscure base two in the width-three example, which tokens must be hidden?
\textbf{Answer:} all three overlapping spans zero, one and two, including cached
features derived from them.
\item Prove the tied-feature RC law.
\textbf{Answer:} complement preserves $g$ and negates $\sigma$; orientation
negates; reversal permutes positions, leaving even and negating odd channels.
\item Must orientation metadata also transform?
\textbf{Answer:} yes for the declared joint action. Keeping $o$ fixed leaves the
learned $os$ contribution with the wrong sign.
\item Does an arbitrary downstream linear map preserve parity?
\textbf{Answer:} not generally; it must commute with the declared sign action,
which forbids unrestricted even/odd mixing.
\item When can invariant pooling be scientifically wrong?
\textbf{Answer:} when the target changes under RC, such as an orientation-specific
label; invariance can erase the required distinction.
\item Compute base centers of $[0,2)$ and $[2,10)$.
\textbf{Answer:} $.5$ and $5.5$, distance five bases.
\item Prove Equation~\ref{evod24:rope}.
\textbf{Answer:} $\mathsf R(a)^\top=\mathsf R(-a)$ and
$\mathsf R(a)\mathsf R(b)=\mathsf R(a+b)$.
\item Is rotary translation invariance an extrapolation benchmark?
\textbf{Answer:} no; it is an algebraic score property, not learned capability at
unseen distance or sequence distributions.
\item Does rotary positioning automatically impose RC equivariance?
\textbf{Answer:} no; reflection negates relative angles and content features also
transform. Their compatible action must be specified.
\item Derive the local-global edge bound.
\textbf{Answer:} local rows contribute at most $T(2w+1)$, global keys at most
$Tg$, and global queries at most $gT$, before overlap removal.
\item Why does a global relay require two layers?
\textbf{Answer:} the first layer computes the global aggregation; only a later
layer can read that new representation.
\item Can causal global edges bypass the triangle?
\textbf{Answer:} no; all edges retain $j\le i$, and any dependency path remains
nonincreasing in base/token time under the legal representation contract.
\item Is a dense masked implementation memory-sparse?
\textbf{Answer:} no; allocating all $T^2$ scores retains dense score storage.
\item What does value/gradient parity establish?
\textbf{Answer:} the sparse and dense references implement the same tested
function and derivatives, not superior speed or genomic prediction.
\item Plan a genomic mechanism ablation.
\textbf{Rubric:} immutable reference/alternate pairs and split groups; explicit
target and strand action; position units; legal information access; conventional
baselines; matched budget and tuning; repeated seeds; attribution/accuracy metrics;
OOD context tests; real memory/latency measurements and retained negative results.
\end{enumerate}

\section{Bridge}
The mechanism contracts now have something concrete to expose: token spans,
orientation action, positions and access edges. Chapter 25 puts those contracts
into a shared experiment envelope while preserving DOGMA's native carry and
Hermon's activation/mask semantics. Fair comparison means the same scientific
question, not erasing each architecture's defining mechanism.
